% !Mode:: "TeX:UTF-8"
\chapter{点云补全算法与三维表面重建算法设计}
\label{cha:algorithm}

% 点云补全作为三维视觉领域的核心任务之一，旨在从部分观测的残缺点云中重建出完整的三维几何结构。随着深度学习技术的快速发展，基于神经网络的点云补全方法取得了显著进展。早期的方法主要依赖于卷积神经网络或图神经网络来提取点云的局部特征，然而这类方法在建模长程依赖关系时存在固有的局限性。Transformer架构凭借自注意力机制的强大表达能力，在自然语言处理和计算机视觉领域取得了突破性成果\cite{vaswani2017attention,dosovitskiy2021vit}。近年来，研究者开始将Transformer引入点云处理领域，并展现出优于传统方法的性能\cite{zhao2021point,guo2021pct}。

% 本章将系统阐述基于Transformer的点云补全算法的整体架构设计，从网络结构、特征编码、渐进式生成策略以及混合损失函数设计四个维度展开深入分析。该算法以AdaPoinTr\cite{yu2023adapointr}为基础框架，在保留原有核心架构的同时，针对损失函数部分进行了系统性优化，旨在提升补全点云的全局结构一致性和局部分布均匀性。

\section{系统架构} % todo:添加示意图

本文的算法设计系统架构可以分为基于Transformer的点云补全模型部分和基于泊松重建的三维表面重建算法部分。其中，基于Transformer的点云补全模型为本文的核心论述重点。本文在PoinTr的基础上，针对点云补全课题中“原始点云残缺区域的补全密度失衡”问题做出了针对性优化，有效解决了在补全过程中出现的点云密度不均匀问题。同时，本文引入了三维表面重建这一点云领域的经典课题作为点云补全的下游任务，旨在直观反映改良后的点云补全算法在实际应用领域中对下游任务的直接提升。

% todo：添加系统处理流程示意图

\subsection{点云补全模型架构}

点云补全任务的形式化定义如下：给定一个部分观测的点云$\mathbf{P}_{partial} = \{\mathbf{p}_i\}_{i=1}^{N_{in}}$，其中$\mathbf{p}_i \in \mathbb{R}^3$表示第$i$个点的三维坐标，$N_{in}$为输入点云的点数，补全网络的目标是学习一个映射函数$f: \mathbb{R}^{N_{in} \times 3} \rightarrow \mathbb{R}^{N_{out} \times 3}$，使得输出点云$\mathbf{P}_{complete} = f(\mathbf{P}_{partial})$尽可能接近真实的完整点云$\mathbf{P}_{gt}$。该任务的核心挑战在于输入点云通常只包含物体表面的一部分，缺失区域可能占据总体积的相当比例，模型需要在缺乏直接观测信息的情况下推断出合理的几何结构。

整体网络架构采用编码器-解码器（Encoder-Decoder）的设计范式，如图\ref{fig:overall_architecture}所示。编码器负责将无序的输入点云转换为具有丰富语义信息的潜在特征表示，解码器则基于该特征表示逐步生成完整的点云。这种架构设计的核心思想在于将复杂的补全任务分解为特征理解和几何生成两个相对独立的子问题，从而降低模型的学习难度。编码器通过多层非线性变换提取输入点云的高层语义特征，这些特征不仅编码了可见区域的几何信息，还隐含了缺失区域的结构线索。解码器则利用这些特征作为条件，通过渐进式生成策略重建完整的三维形状。
\FloatBarrier
\begin{figure}[htbp]
\centering
\resizebox{0.85\textwidth}{!}{%
\begin{tikzpicture}[
    >=stealth,
    node distance=2.5cm,
    block/.style={
        rectangle,
        draw=black!70,
        fill=blue!10,
        minimum width=3.5cm,
        minimum height=1.2cm,
        rounded corners=2pt,
        font=\small,
        align=center
    },
    process/.style={
        rectangle,
        draw=black!70,
        fill=green!10,
        minimum width=3.5cm,
        minimum height=1.2cm,
        rounded corners=2pt,
        font=\small,
        align=center
    },
    output/.style={
        rectangle,
        draw=black!70,
        fill=orange!15,
        minimum width=3.5cm,
        minimum height=1.2cm,
        rounded corners=2pt,
        font=\small,
        align=center
    },
    arrow/.style={
        ->,
        thick,
        >=stealth
    }
]

% 输入层
\node[block] (input) {输入点云 $\mathbf{P}_{partial}$\\$(N_{in}=2048, 3)$};

% 编码器部分
\node[block, below=1.2cm of input] (fps) {最远点采样 FPS};

\node[block, below=1.2cm of fps] (knn) {K近邻分组};

\node[block, below=1.2cm of knn] (patch) {Patch特征编码};

\node[block, below=1.2cm of patch] (pos) {三维位置编码};

\node[process, below=1.2cm of pos] (encoder) {几何感知Transformer编码器\\$(6 \text{ 层}, \text{多头注意力 } h=6)$};

% 全局特征
\node[process, below=1.2cm of encoder] (maxpool) {最大池化 Max Pooling\\全局特征 $\mathbf{F}_{global} \in \mathbb{R}^{d}$};

% 粗粒度生成
\node[output, below=1.5cm of maxpool] (coarse) {粗粒度点云生成\\$(N_{coarse}=512, 3)$};

% 自适应查询生成
\node[process, below=1.2cm of coarse, fill=purple!10] (query) {自适应查询生成\\查询排序 + 去噪任务};

% 细粒度生成
\node[output, below=1.2cm of query] (decoder) {Transformer解码器\\$(8 \text{ 层自注意力+交叉注意力})$};

\node[output, below=1.2cm of decoder] (fine) {细粒度点云输出\\$(N_{fine}=8192, 3)$};

% 损失函数
\node[block, right=2.5cm of coarse, fill=red!10] (loss_coarse) {粗粒度损失\\$\mathcal{L}_{coarse} = \lambda_1 \mathcal{L}_{EMD} + \lambda_2 \mathcal{L}_{CD}$};

\node[block, right=2.5cm of query, fill=red!10] (loss_denoise) {去噪损失\\$\mathcal{L}_{denoise}$};

\node[block, right=2.5cm of fine, fill=red!10] (loss_fine) {细粒度损失\\$\mathcal{L}_{fine} = \lambda_3 \mathcal{L}_{DCD} + \lambda_4 \mathcal{L}_{CD} + \lambda_5 \mathcal{L}_{rep}$};

% 箭头连接
\draw[arrow] (input) -- (fps);
\draw[arrow] (fps) -- (knn);
\draw[arrow] (knn) -- (patch);
\draw[arrow] (patch) -- (pos);
\draw[arrow] (pos) -- (encoder);
\draw[arrow] (encoder) -- (maxpool);
\draw[arrow] (maxpool) -- (coarse);
\draw[arrow] (coarse) -- (query);
\draw[arrow] (query) -- (decoder);
\draw[arrow] (decoder) -- (fine);

% 损失函数连接
\draw[arrow, dashed] (coarse) -- (loss_coarse);
\draw[arrow, dashed] (query) -- (loss_denoise);
\draw[arrow, dashed] (fine) -- (loss_fine);

% 标注
\node[above=0.3cm of input, font=\small\bfseries] {输入};
\node[left=0.5cm of encoder, font=\small\bfseries, text=blue!70!black] {编码器};
\node[left=0.5cm of query, font=\small\bfseries, text=purple!70!black] {查询优化};
\node[left=0.5cm of coarse, font=\small\bfseries, text=orange!70!black] {解码器};
\node[right=0.3cm of loss_coarse, font=\small, text=red!70!black] {训练阶段};
\node[right=0.3cm of loss_denoise, font=\small, text=red!70!black] {训练阶段};
\node[right=0.3cm of loss_fine, font=\small, text=red!70!black] {训练阶段};

\end{tikzpicture}%
}
\caption{基于Transformer的点云补全网络整体架构图}
\label{fig:overall_architecture}
\end{figure}
\FloatBarrier

编码器的输入为$N_{in}=2048$个点的残缺点云，经过多层几何感知Transformer模块的处理后，输出一个全局特征向量$\mathbf{F}_{global} \in \mathbb{R}^{d}$，其中$d$为特征维度。该特征向量编码了输入点云的全局几何结构和语义信息，为后续的解码过程提供必要的上下文指导。在特征提取过程中，每个Transformer层通过自注意力机制建立点与点之间的长程依赖关系，使得模型能够捕获跨越整个点云的几何关联。这种全局建模能力是传统卷积方法所欠缺的，因为卷积操作的感受野受限于卷积核的大小，需要通过多层堆叠才能逐步扩大感受野，而自注意力机制在单层中即可建立任意两点之间的直接联系。

解码器采用Coarse-to-Fine的渐进式生成策略，分为粗粒度生成和细粒度优化两个阶段。粗粒度阶段生成$N_{coarse}=512$个点的骨架点云$\mathbf{P}_{coarse}$，这些点构成了补全物体的基本形状轮廓。细粒度阶段则在粗粒度点云的基础上进行上采样和细节优化，最终输出$N_{fine}=8192$个点的完整点云$\mathbf{P}_{fine}$。这种由粗到精的生成方式使得模型能够先关注全局结构的正确性，再逐步完善局部细节，有效避免了直接生成大量点云时可能出现的结构混乱问题。从优化的角度来看，粗粒度阶段的损失函数主要关注全局拓扑对齐，而细粒度阶段则更加注重局部细节的精确性和点云分布的均匀性。

在训练阶段，模型通过混合损失函数进行优化。粗粒度点云与真实点云之间计算EMD（Earth Mover's Distance）损失，确保骨架点云的全局拓扑对齐；细粒度点云则同时计算DCD（Density-aware Chamfer Distance）损失、标准CD损失和排斥损失（Repulsion Loss），以兼顾补全精度和点云分布的均匀性。多损失协同优化的策略使得模型能够在不同粒度层面上获得有效的梯度信号，从而提升整体的补全质量。这些优化措施旨在降低补全点云的Chamfer Distance指标，提升补全结果的视觉质量和几何精度。

\subsection{三维表面重建算法架构}

表面重建算法的整体架构可划分为三个层次：数据预处理、核心重建与后处理优化。图\ref{fig:reconstruction_pipeline}展示了从补全点云到最终网格模型的完整处理流程。
\FloatBarrier
\begin{figure}[htbp]
    \centering
    \resizebox{0.9\textwidth}{!}{%
    \begin{tikzpicture}[
        node distance=0.8cm,
        box/.style={
            rectangle,
            rounded corners,
            draw=black,
            thick,
            minimum width=2.8cm,
            minimum height=0.9cm,
            align=center,
            fill=#1,
            font=\small
        },
        box/.default=blue!10,
        arrow/.style={
            ->,
            >=stealth,
            thick
        }
    ]
    % 第一层：数据预处理
    \node[box=blue!10] (input) {补全点云输入\\8192点};
    \node[box=blue!10, right=1.5cm of input] (normal) {法线估计与定向};
    
    % 第二层：核心重建
    \node[box=green!10, below=1cm of normal] (poisson) {泊松表面重建};
    
    % 第三层：后处理优化
    \node[box=orange!10, below=1cm of poisson] (density) {密度阈值过滤};
    \node[box=orange!10, right=1.5cm of density] (smooth) {Taubin平滑};
    \node[box=orange!10, right=1.5cm of smooth] (output) {网格输出};
    
    % 箭头连接
    \draw[arrow] (input) -- (normal);
    \draw[arrow] (normal) -- (poisson);
    \draw[arrow] (poisson) -- (density);
    \draw[arrow] (density) -- (smooth);
    \draw[arrow] (smooth) -- (output);
    
    % 左侧标注
    \node[left=0.3cm of input, text width=1.8cm, align=center, font=\small] (label1) {数据\\预处理};
    \draw[dashed, gray] (label1.north east) -- (input.north west);
    \draw[dashed, gray] (label1.south east) -- (input.south west);
    
    \node[left=0.3cm of poisson, text width=1.8cm, align=center, font=\small] (label2) {核心\\重建};
    \draw[dashed, gray] (label2.north east) -- (poisson.north west);
    \draw[dashed, gray] (label2.south east) -- (poisson.south west);
    
    \node[left=0.3cm of density, text width=1.8cm, align=center, font=\small] (label3) {后处理\\优化};
    \draw[dashed, gray] (label3.north east) -- (density.north west);
    \draw[dashed, gray] (label3.south east) -- (density.south west);
    \end{tikzpicture}%
    }
    \caption{三维表面重建算法处理流程}
    \label{fig:reconstruction_pipeline}
\end{figure}
\FloatBarrier

算法流程从补全模型输出的8192个点云数据开始，依次经过法线估计与定向、泊松表面重建、密度阈值过滤和Taubin平滑四个核心步骤，最终输出三角网格模型。各步骤的具体原理与参数配置将在后续小节中详细展开。

\section{基于 Transformer 的点云补全算法设计}

\subsection{点云序列的Transformer编码}

点云数据与图像或文本等规则数据结构存在本质差异。点云是一种无序的集合，其点的排列顺序不具有语义意义，同时点的空间分布往往呈现非均匀的特性。这些特性使得直接将标准Transformer应用于点云处理面临诸多挑战。点云序列的Transformer编码模块通过序列化处理和几何感知注意力机制，有效地解决了这些问题，为后续的补全过程提供了高质量的特征表示。

(1) 点云序列化与位置编码

点云序列化的核心目标是将无序的点集转换为有序的序列，同时保留点云的局部几何结构信息。这一过程借鉴了Vision Transformer\cite{dosovitskiy2021vit}中将图像分割为patch的思想，但需要根据点云的特性进行适应性改造。图像数据具有天然的网格结构，可以均匀地划分为固定大小的patch，而点云数据缺乏这种规则的空间组织方式，因此需要设计专门的序列化策略来捕获点云的局部几何特征。

具体而言，序列化过程首先通过最远点采样（Farthest Point Sampling, FPS）从输入点云中选择$M$个中心点。FPS算法通过迭代选择距离已选点集最远的点，确保中心点在空间中的分布尽可能均匀。与随机采样相比，FPS能够更好地覆盖点云的整体空间范围，避免采样点在某些区域过度集中而在其他区域完全缺失。对于每个中心点，通过K近邻搜索获取其周围的$K$个邻近点，形成一个局部点云片（patch）。这种基于空间邻域关系的分组方式使得每个patch能够捕获局部的几何结构信息，例如曲率变化、表面法线方向以及局部形状特征。

每个patch中的点通过一个共享权重的多层感知机（MLP）进行特征提取，将$K$个三维点映射为一个$d$维的特征向量。这一过程类似于PointNet\cite{qi2017pointnet}中的点特征编码，但作用于局部patch而非全局点云。共享权重的设计确保了无论patch位于点云的哪个位置，特征提取的方式保持一致，从而保证了特征的平移不变性。经过特征提取后，原始的$N_{in}$个点被转换为$M$个patch特征，形成了长度为$M$的特征序列。这个序列作为Transformer编码器的输入，每个patch特征对应序列中的一个token。

位置编码是Transformer架构中不可或缺的组成部分。由于自注意力机制本身不具备感知输入元素位置关系的能力，必须通过显式的位置编码将空间信息注入特征序列。对于自然语言处理任务，位置编码通常采用正弦余弦函数或可学习的一维位置索引，因为文本数据具有天然的线性顺序。对于点云数据，位置编码的设计需要反映点的三维空间坐标。本算法采用可学习的三维位置编码方案，将每个patch中心点的三维坐标通过一个MLP映射为与特征维度相同的编码向量，然后与patch特征进行逐元素相加。这种位置编码方式使得模型能够在自注意力计算中感知patch之间的空间相对位置，从而更好地建模点云的几何结构。

位置编码的设计对模型的性能具有重要影响。如果位置编码过于简单，例如仅使用一维索引，则无法充分表达点云的三维空间关系；如果位置编码过于复杂，例如直接输入原始坐标而不进行非线性变换，则可能引入不必要的噪声。可学习的三维位置编码方案通过在训练过程中自适应地调整编码参数，能够在保留空间信息的同时过滤掉对任务无关的坐标细节。此外，位置编码与patch特征的逐元素相加操作确保了空间信息和语义信息在特征空间中的融合，使得后续的自注意力计算能够同时考虑两者的影响。

(2) 几何感知自注意力模块

标准Transformer的自注意力机制通过计算查询（Query）和键（Key）之间的点积来衡量元素之间的相关性，这种相关性度量主要基于特征的语义相似度，而忽略了元素之间的空间几何关系。对于点云补全任务而言，空间位置相近的点往往具有更强的几何关联，例如同一物体表面的相邻区域通常具有相似的曲率和法线方向。因此在注意力计算中引入几何感知能力是提升模型性能的关键。

几何感知自注意力模块在标准自注意力的基础上增加了空间位置偏置项。具体而言，对于第$i$个patch和第$j$个patch之间的注意力权重计算如下：
\begin{equation}
\alpha_{ij} = \text{softmax}_j\left(\frac{\mathbf{q}_i^\top \mathbf{k}_j + \phi(\mathbf{c}_i, \mathbf{c}_j)}{\sqrt{d}}\right)
\end{equation}

其中$\mathbf{q}_i = \mathbf{W}_Q \mathbf{f}_i$和$\mathbf{k}_j = \mathbf{W}_K \mathbf{f}_j$分别为查询和键向量，$\mathbf{f}_i$和$\mathbf{f}_j$为patch特征，$\mathbf{c}_i$和$\mathbf{c}_j$为patch中心点的三维坐标，$\phi(\cdot, \cdot)$为位置偏置函数，通常实现为一个MLP，将两个中心点的相对坐标映射为一个标量偏置值。

位置偏置项$\phi(\mathbf{c}_i, \mathbf{c}_j)$的引入使得注意力机制能够同时考虑特征相似度和空间邻近度。当两个patch在空间中距离较近时，位置偏置项会增大它们之间的注意力权重，促使模型更多地关注局部几何结构；当两个patch距离较远时，位置偏置项的影响减弱，注意力权重主要由特征相似度决定，从而保留了建模长程依赖的能力。这种设计使得模型能够在局部细节和全局结构之间取得平衡，既不会过度关注局部区域而忽略整体形状，也不会因为长程依赖的建模而丢失局部几何信息。

此外，几何感知自注意力模块还采用了多头注意力（Multi-Head Attention）机制，将特征和位置信息投影到多个子空间中并行计算注意力，然后将结果拼接并通过线性变换进行融合。多头机制使得模型能够从不同的角度捕获点云的几何特征，例如一个注意力头可能专注于局部表面法线的变化，另一个头可能关注全局形状的对称性，还有一个头可能负责捕获物体部件之间的空间关系。这种多角度的特征提取增强了模型的表达能力，使得编码器能够生成更加丰富和 discriminative 的特征表示。

自注意力机制的计算复杂度为$O(M^2 d)$，其中$M$为patch数量，$d$为特征维度。对于$M=256$的典型配置，计算开销在可接受范围内。然而当patch数量增加到数千级别时，二次复杂度的计算将成为性能瓶颈。本算法通过合理选择patch数量，在保证特征表达能力的同时控制了计算开销。同时，多头注意力的并行计算特性使得该模块能够充分利用现代GPU的并行计算能力，进一步提升了计算效率。

(3) 全局特征提取

经过多层几何感知Transformer模块的处理后，每个patch的特征已经融合了来自其他patch的信息，形成了具有全局上下文感知的特征表示。为了获得整个点云的全局特征，需要对这些patch特征进行聚合。全局特征提取的目标是将$M$个patch特征压缩为一个固定维度的向量，该向量需要保留输入点云的关键语义信息和几何结构线索。

本算法采用最大池化（Max Pooling）操作进行特征聚合。最大池化通过对每个特征维度取所有patch特征的最大值，得到一个$d$维的全局特征向量。这种聚合方式具有排列不变性，即无论patch的排列顺序如何，聚合结果保持一致，这与点云的无序性特性相一致。同时，最大池化能够保留最显著的特征响应，使得全局特征对局部几何变化具有一定的鲁棒性。与平均池化相比，最大池化更倾向于保留那些在某个patch中表现突出的特征，这些特征通常对应于点云中的关键几何结构，例如尖锐的边缘、突出的角点或独特的表面纹理。

除了最大池化之外，本算法还保留了每个patch的特征作为局部特征表示，用于后续解码器中的细粒度生成。全局特征和局部特征的结合使得解码器能够在不同尺度上利用编码器的输出，从而实现更精确的点云补全。全局特征提供了物体整体形状的宏观指导，而局部特征则保留了各个区域的微观几何细节，两者的协同作用使得解码器能够在生成过程中同时考虑全局一致性和局部精确性。

编码器输出的全局特征向量$\mathbf{F}_{global}$不仅包含了输入点云的几何结构信息，还编码了缺失区域的语义线索。这些线索对于指导解码器生成合理的补全结果至关重要。例如，当输入点云显示物体的左半部分时，全局特征中会隐含右半部分的对称性信息，解码器可以利用这些信息推断出缺失区域的形状。对于具有对称结构的物体（如飞机、汽车、椅子等），这种对称性线索尤为有效；对于非对称物体，全局特征中编码的类别语义信息也能帮助解码器生成符合该类别典型形状的补全结果。

多层Transformer模块的堆叠进一步增强了特征提取的能力。每一层Transformer模块通过自注意力机制和多层感知机对特征进行非线性变换，使得高层特征能够捕获更加抽象和语义化的信息。随着层数的增加，特征的感知范围逐步扩大，从局部几何细节逐渐过渡到全局形状结构。这种层次化的特征提取方式与人类视觉系统从边缘检测到形状识别的处理过程具有一定的相似性，体现了深度学习模型在特征学习方面的强大能力。

\subsection{Coarse-to-Fine渐进式生成策略}

直接生成高密度的完整点云面临诸多挑战。当输出点数达到数千甚至上万时，模型需要同时考虑全局结构的正确性和局部细节的精确性，这要求模型具备极强的表达能力和泛化能力。Coarse-to-Fine渐进式生成策略通过将补全过程分解为多个阶段，逐步增加点云的密度和细节，有效降低了每个阶段的学习难度。这种策略的思想源于人类认知过程中从整体到局部的观察方式，也符合三维形状生成的自然规律。从优化的角度来看，渐进式生成策略使得损失函数能够在不同粒度层面上提供有针对性的梯度信号，避免了单一尺度下梯度方向可能存在的冲突。

(1) 粗粒度形状生成

粗粒度形状生成阶段的目标是从全局特征$\mathbf{F}_{global}$中生成$N_{coarse}=512$个点的骨架点云$\mathbf{P}_{coarse}$。这些点构成了补全物体的基本形状轮廓，决定了补全结果的全局拓扑结构。粗粒度点云的质量直接影响后续细粒度优化的效果。如果粗粒度点云的全局结构存在偏差，细粒度阶段很难通过局部调整来纠正这些偏差。因此，粗粒度生成阶段需要足够强的约束来确保生成的骨架点云与真实形状在全局拓扑上保持一致。

具体而言，粗粒度生成器首先将全局特征向量$\mathbf{F}_{global}$通过一个MLP扩展为一个种子点集合。种子点的数量通常远少于最终输出的粗粒度点数，例如64个种子点。每个种子点携带了全局特征的部分信息，代表了物体形状的不同区域。种子点的生成过程可以视为对物体形状的粗略划分，每个种子点对应于物体的一个主要部件或区域。随后，种子点通过特征折叠（Folding）操作进行扩展。特征折叠借鉴了二维流形展开的思想，将每个种子点与一组二维网格坐标拼接，通过解码器网络映射为三维空间中的多个点。这种操作使得每个种子点能够生成一个局部的点云片，所有种子点生成的点云片组合在一起形成了完整的粗粒度点云。

特征折叠操作的数学本质是学习一个从二维参数空间到三维物理空间的映射函数。二维网格坐标提供了局部点云片的参数化表示，解码器网络则学习如何将这些参数化坐标转换为实际的三维点坐标。这种设计使得模型能够利用二维流形的连续性来保证生成的点云片在三维空间中具有合理的几何结构。与直接回归点坐标相比，特征折叠能够更好地保持局部点云的拓扑一致性，减少点的交叉和重叠现象。

粗粒度点云的生成过程可以形式化地表示为：
\begin{equation}
\mathbf{P}_{coarse} = \text{Decoder}_{coarse}(\mathbf{F}_{global}, \mathbf{G}_{2D})
\end{equation}

其中$\mathbf{G}_{2D} \in \mathbb{R}^{N_{coarse} \times 2}$为二维网格坐标，$\text{Decoder}_{coarse}$为粗粒度解码器网络，通常由多层MLP组成。

在粗粒度生成过程中，全局特征$\mathbf{F}_{global}$作为条件输入，指导种子点的生成和特征折叠的方向。这种条件生成方式使得粗粒度点云能够与输入点云的可见区域保持几何一致性，同时根据缺失区域的语义线索推断出合理的补全形状。对于具有对称结构的物体，全局特征中编码的对称性信息能够引导模型生成对称的补全结果；对于具有重复部件的物体（如椅子的四条腿、桌子的多个支撑柱等），全局特征中的重复模式信息能够帮助模型生成结构一致的多个部件。

(2) 自适应查询生成机制

在粗粒度点云生成过程中，查询点（Query Points）的选择策略对补全质量具有决定性影响。早期的点云补全方法通常采用固定的查询生成策略，例如通过最远点采样（FPS）从输入点云中均匀采样一定数量的点作为查询点。这种固定策略虽然简单高效，但缺乏对不同几何区域的自适应能力，无法根据点云的局部复杂度动态调整查询点的分布。本文受AdaPoinTr\cite{yu2023adapointr}启发，设计了自适应查询生成机制，旨在提升模型的整体补全质量与鲁棒性。

自适应查询排序机制通过重要性评估来优化查询点的选择。具体而言，模型首先生成初始的粗糙点云，其点数通常多于最终用于解码的查询点数量。随后，使用一个轻量级的排序网络对每个粗糙点进行重要性评分：
\begin{equation}
s_i = \sigma(\text{MLP}_{rank}(\mathbf{p}_i^{coarse}))
\end{equation}

其中$\mathbf{p}_i^{coarse}$为第$i$个粗糙点的坐标，$\text{MLP}_{rank}$为多层感知机，$\sigma(\cdot)$为Sigmoid激活函数，$s_i \in (0, 1)$为该点的重要性分数。排序网络根据重要性分数对所有粗糙点进行降序排列，选择排名靠前的$N_{query}$个点作为最终的查询点集。

这种自适应排序机制的核心优势在于能够优先选择几何特征丰富区域的点作为查询。例如，在物体的边缘、角点或曲率变化剧烈的区域，这些点携带了更多的形状信息，对于指导后续的细粒度生成具有更高的价值。通过排序网络的学习，模型能够自动识别这些关键区域，并将有限的查询点资源分配给最具信息量的位置。与固定FPS采样相比，自适应查询排序使得查询点的分布更加符合物体的几何复杂度分布，从而提升了整体补全质量。

去噪任务辅助训练是自适应查询生成机制的另一项关键创新。该机制借鉴了自然语言处理中降噪自动编码器（Denoising Autoencoder）的思想，通过在训练过程中引入额外的去噪目标来增强模型的鲁棒性和泛化能力。具体实现方式为：在训练阶段，从真实完整点云中通过FPS采样$N_{noise}$个点，并对这些点添加随机噪声扰动：
\begin{equation}
\mathbf{p}_i^{noise} = \mathbf{p}_i^{gt} + \epsilon_i, \quad \epsilon_i \sim \mathcal{N}(0, \sigma^2)
\end{equation}

其中$\mathbf{p}_i^{gt}$为真实点，$\epsilon_i$为服从正态分布的噪声，$\sigma$为噪声强度控制参数。这些带噪声的点与模型生成的粗糙点云拼接在一起，共同作为解码器的输入。解码器在生成细粒度点云的过程中，需要同时完成两个任务：从粗糙点生成完整点云（重建任务）和从噪声点恢复正确位置（去噪任务）。

去噪任务的引入为模型提供了额外的监督信号。在训练过程中，模型不仅需要学习如何从部分观测推断完整形状，还需要学习如何从带噪声的输入中恢复准确的几何位置。这种多任务学习策略使得模型能够捕获更加鲁棒的几何特征表示，增强了对输入扰动的容忍度。

自适应查询排序和去噪任务辅助训练之间存在协同效应。自适应查询排序确保了查询点的质量，使得解码器能够基于最具信息量的查询点进行细粒度生成；去噪任务则通过额外的监督信号增强了模型对查询点特征的利用效率。两者的结合使得模型能够在保持较高计算效率的同时，提高整体补全质量与鲁棒性。

(3) 细粒度细节优化

细粒度细节优化阶段在粗粒度点云$\mathbf{P}_{coarse}$的基础上进行上采样和细节优化，最终输出$N_{fine}=8192$个点的完整点云$\mathbf{P}_{fine}$。这一阶段的核心挑战在于如何在增加点云密度的同时，保持点云分布的均匀性和表面细节的精确性。与粗粒度阶段关注全局拓扑不同，细粒度阶段需要处理更加精细的几何特征，例如表面曲率的连续变化、边缘的锐利程度以及局部纹理的精确重建。

细粒度优化器首先通过K近邻搜索为每个粗粒度点获取其局部邻域信息。对于每个粗粒度点$\mathbf{p}_i^{coarse}$，搜索其周围的$K$个最近邻点，并将这些邻域点的特征与全局特征进行拼接，形成增强的局部特征表示。这种局部特征增强方式使得模型能够在生成细粒度点时充分考虑周围的几何上下文。邻域信息的引入使得每个细粒度点的生成不仅依赖于其对应的粗粒度点，还受到周围点的影响，从而保证了生成点云的局部连续性和光滑性。

随后，每个粗粒度点通过上采样操作生成多个细粒度点。上采样操作借鉴了PU-Net\cite{yu2018punet}中的特征扩展思想，通过多层MLP将单个粗粒度点的特征扩展为$r$个细粒度点的特征，其中$r = N_{fine} / N_{coarse}$为上采样倍率。扩展后的特征通过坐标重建模块转换为三维空间中的点坐标。上采样倍率的选择需要在补全精度和计算开销之间取得平衡。过高的倍率会导致单个粗粒度点需要生成过多的细粒度点，增加了局部重建的难度；过低的倍率则需要更多的粗粒度点作为基础，增加了粗粒度阶段的计算负担。

细粒度点的生成过程可以表示为：
\begin{equation}
\mathbf{P}_{fine} = \text{Decoder}_{fine}(\mathbf{P}_{coarse}, \mathbf{F}_{local}, \mathbf{F}_{global})
\end{equation}

其中$\mathbf{F}_{local}$为局部特征表示，$\text{Decoder}_{fine}$为细粒度解码器网络。

在细粒度优化过程中，点云分布的均匀性是一个需要特别关注的问题。由于Transformer模型的注意力机制倾向于在某些区域过度集中，生成的点云容易出现局部堆积和空洞并存的现象。这种不均匀分布不仅影响补全结果的视觉质量，还会导致下游任务（如分类、分割、配准等）的性能下降。为了解决这一问题，本算法在细粒度阶段引入了排斥损失（Repulsion Loss），通过惩罚过近的点对来促使点云均匀分布。排斥损失的作用机制类似于物理系统中的排斥力，当两个点之间的距离小于某个阈值时，会产生一个排斥力将它们推开，从而避免点的过度聚集。

密度感知倒角距离（DCD）损失也在细粒度阶段发挥重要作用。DCD通过密度感知权重防止多个预测点匹配到同一个真实点，从而进一步提升点云分布的质量。与标准CD相比，DCD对局部密度变化更加敏感，能够有效识别和惩罚那些导致局部点云堆积的预测结果。排斥损失和DCD损失的协同作用使得细粒度点云在全局结构正确的基础上，具备了更加均匀和自然的点分布特性。

Coarse-to-Fine渐进式生成策略的优势在于将复杂的补全任务分解为多个相对简单的子任务。粗粒度阶段负责建立正确的全局拓扑，细粒度阶段负责完善局部细节和点分布。这种分解使得每个子任务的优化目标更加明确，梯度信号更加稳定，从而提升了模型的整体训练效率和补全质量。同时，渐进式生成策略也为混合损失函数的设计提供了自然的基础，使得不同粒度层面可以采用不同的损失组合，实现更加精细化的优化控制。

\subsection{混合损失函数设计}

损失函数的设计是点云补全模型训练的核心环节。合适的损失函数不仅能够引导模型学习到正确的补全策略，还能够影响生成点云的分布特性和细节质量。传统的点云补全方法主要依赖于Chamfer Distance（CD）作为损失函数，然而CD存在对局部密度不敏感、容易产生多点匹配同一真实点等缺陷\cite{wu2021densityaware}。为了克服这些局限性，本算法设计了混合损失函数，在不同粒度层面上采用不同的损失组合，并通过分阶段权重调度实现多损失的协同优化。

点云补全任务的损失函数设计需要考虑多个维度的约束。从全局层面来看，补全点云需要与真实点云在整体形状和拓扑结构上保持一致；从局部层面来看，补全点云的点分布需要均匀，不能出现局部堆积或空洞；从细节层面来看，补全点云需要精确地重建物体表面的几何特征，例如边缘、角点和曲面过渡区域。单一的损失函数往往只能关注其中某一个或某几个维度，难以同时满足所有约束条件。混合损失函数通过组合多种不同的损失项，能够在多个维度上同时施加约束，从而获得更全面的优化效果。

(1) 粗粒度损失函数设计

粗粒度点云$\mathbf{P}_{coarse}$包含512个点，这些点构成了补全物体的基本骨架。对于骨架点云而言，全局拓扑结构的正确性比局部细节的精确性更为重要。因此，粗粒度损失函数的设计重点在于确保生成的骨架点云与真实点云在全局分布上保持一致。

Earth Mover's Distance（EMD）是一种衡量两个点集之间分布差异的度量，其定义为将一个点集变换为另一个点集所需的最小移动代价。与CD不同，EMD要求两个点集之间建立双射匹配关系，即每个预测点必须唯一地匹配到一个真实点。这种双射约束使得EMD对全局分布的敏感性远高于CD，能够有效防止局部点云堆积现象。从最优传输理论的角度来看，EMD等价于求解一个最优传输问题，其中源分布为预测点云，目标分布为真实点云，传输代价为点对之间的欧氏距离。由于EMD的数学定义已在第二章中展开论述，在此不做过多赘述。

% EMD的数学定义为：

% \begin{equation}
% \text{EMD}(\mathbf{P}_{coarse}, \mathbf{P}_{gt}) = \min_{\phi: \mathbf{P}_{coarse} \rightarrow \mathbf{P}_{gt}} \sum_{\mathbf{x} \in \mathbf{P}_{coarse}} \|\mathbf{x} - \phi(\mathbf{x})\|_2
% \end{equation}

% 其中$\phi$为$\mathbf{P}_{coarse}$到$\mathbf{P}_{gt}$的双射映射。

EMD的计算通常采用匈牙利算法（Hungarian Algorithm）或拍卖算法（Auction Algorithm）来求解最优匹配。这些算法的时间复杂度为$O(N^3)$，其中$N$为点集的点数。对于粗粒度点云的512个点，计算开销在可接受范围内。然而当点数增加到数千级别时，EMD的计算将成为性能瓶颈。这也是为什么EMD通常只适用于粗粒度阶段，而不适合直接用于细粒度点云的损失计算。

然而，EMD的计算要求两个点集的点数相等。在粗粒度阶段，$\mathbf{P}_{coarse}$包含512个点，而真实点云$\mathbf{P}_{gt}$包含8192个点。为了应用EMD损失，需要从真实点云中采样出512个点作为匹配目标。传统的随机采样方法虽然简单，但采样结果的随机性较大，可能导致训练过程中的梯度不稳定。本算法采用最远点采样（FPS）从真实点云中选取512个点，FPS能够确保采样点在空间中均匀分布，从而提供更稳定的匹配目标。与随机采样相比，FPS采样的点集具有更好的空间覆盖性，能够更全面地代表真实点云的整体形状，使得EMD损失能够更准确地评估粗粒度点云的全局拓扑质量。

粗粒度损失函数由EMD损失和CD损失加权组合而成：
\begin{equation}
\mathcal{L}_{coarse} = w_{emd} \cdot \text{EMD}(\mathbf{P}_{coarse}, \text{FPS}(\mathbf{P}_{gt})) + w_{cd} \cdot \text{CD}(\mathbf{P}_{coarse}, \mathbf{P}_{gt})
\end{equation}

其中$w_{emd}$和$w_{cd}$分别为EMD和CD的权重系数，$\text{FPS}(\mathbf{P}_{gt})$表示从真实点云中通过FPS采样得到的512个点。

EMD和CD的组合使用体现了全局约束和局部约束的互补性。EMD通过双射匹配确保全局分布的一致性，但对局部细节的敏感性较弱；CD通过最近邻匹配捕捉局部几何特征，但对全局分布的约束力不足。两者的结合使得粗粒度损失函数能够同时关注全局拓扑和局部几何，为粗粒度点云的生成提供更全面的优化指导。

在训练的不同阶段，EMD和CD的权重会进行动态调整。在训练初期，给予EMD更高的权重，帮助模型快速建立正确的全局拓扑结构；随着训练的推进，逐渐降低EMD的权重，使模型能够更多地关注局部细节的优化。这种分阶段权重调度策略使得模型能够在不同训练阶段获得最合适的优化目标。训练初期的高EMD权重促使编码器快速学习到有效的全局特征表示，而后期的高CD权重则帮助粗粒度解码器完善局部几何细节。

(2) 细粒度损失函数设计

细粒度点云$\mathbf{P}_{fine}$包含8192个点，是最终的补全输出。对于细粒度点云而言，不仅需要保证全局结构的正确性，还需要确保点云分布的均匀性和局部细节的精确性。因此，细粒度损失函数的设计需要综合考虑多个方面的约束。

密度感知倒角距离（Density-aware Chamfer Distance, DCD）是细粒度损失函数的核心组成部分。DCD由Wu等人\cite{wu2021densityaware}提出，旨在解决标准CD对局部密度不敏感的问题。标准CD的计算方式为预测点到真实点的最近距离之和加上真实点到预测点的最近距离之和，这种计算方式允许多个预测点匹配到同一个真实点，从而导致局部点云堆积。DCD通过引入密度感知权重来解决这一问题。

原始DCD的数学定义为：
\begin{equation}
\text{DCD}(\mathbf{P}_{fine}, \mathbf{P}_{gt}) = \frac{1}{N}\sum_{i=1}^{N}\frac{1}{\sum_{j=1}^{M}\exp(-\gamma\|\mathbf{x}_i-\mathbf{y}_j\|^2)} + \frac{1}{M}\sum_{j=1}^{M}\frac{1}{\sum_{i=1}^{N}\exp(-\gamma\|\mathbf{y}_j-\mathbf{x}_i\|^2)}
\end{equation}

其中$\mathbf{x}_i \in \mathbf{P}_{fine}$，$\mathbf{y}_j \in \mathbf{P}_{gt}$，$\gamma$为温度参数，控制指数衰减的速度。

原始DCD的物理意义可以从概率分布的角度来理解。指数项$\exp(-\gamma\|\mathbf{x}_i-\mathbf{y}_j\|^2)$可以视为一个以$\mathbf{x}_i$为中心、方差由$\gamma$控制的高斯（Gaussian）核函数。对于每个预测点$\mathbf{x}_i$，所有真实点$\mathbf{y}_j$对该点的高斯核响应之和反映了$\mathbf{x}_i$周围真实点的密度分布。如果多个预测点都靠近同一个真实点，则该真实点周围的预测点密度较高，对应的高斯核响应之和较大，取倒数后得到的密度感知权重较小，从而降低了这些预测点对总损失的贡献。这种机制有效地惩罚了局部点云堆积现象，促使预测点更加均匀地分布在真实点云的表面。

原始DCD的计算复杂度为$O(N \times M)$，对于$N=M=8192$且批量大小为64的情况，需要创建$64 \times 8192 \times 8192$的距离矩阵，这将导致严重的显存溢出问题。为了在保持DCD核心思想的同时降低计算复杂度，本算法提出了一种基于KNN的简化版DCD实现。

简化版DCD的核心思想是：当温度参数$\gamma$足够大时，指数项$\exp(-\gamma\|\mathbf{x}_i-\mathbf{y}_j\|^2)$对于距离较远的点对会快速衰减到接近零，因此只有最近邻的点对指数和有显著贡献。基于这一观察，简化版DCD仅使用最近邻匹配来计算密度感知权重，将计算复杂度从$O(N \times M)$降低到$O(N+M)$。

具体而言，简化版DCD首先通过KNN搜索（$K=1$）获取每个预测点的最近邻真实点及其距离，同时获取每个真实点的最近邻预测点及其距离。然后统计每个真实点被预测点引用的次数，引用次数越多的点说明其周围预测点密度越高，对应的密度感知权重越低。最终的损失计算为：
\begin{equation}
\mathcal{L}_{dcd} = \frac{1}{N}\sum_{i=1}^{N}\left(1 - \exp(-\gamma\|\mathbf{x}_i-\mathbf{y}_{k(i)}\|^2) \cdot w_i\right) + \frac{1}{M}\sum_{j=1}^{M}\left(1 - \exp(-\gamma\|\mathbf{y}_j-\mathbf{x}_{k(j)}\|^2) \cdot w_j\right)
\end{equation}

其中$\mathbf{y}_{k(i)}$为$\mathbf{x}_i$的最近邻真实点，$\mathbf{x}_{k(j)}$为$\mathbf{y}_j$的最近邻预测点，$w_i$和$w_j$为密度感知权重，定义为引用次数的倒数。

温度参数$\gamma$的取值对DCD的性能有重要影响。根据Wu等人\cite{wu2021densityaware}的分析，对于ShapeNet等归一化数据集，最近邻点对的距离平方通常在$10^{-4}$到$10^{-3}$之间，因此设置$\gamma=1000$能够将指数项映射到$[0.37, 0.90]$的合理区间，既保证了足够的区分度，又避免了梯度消失问题。本算法在实现中严格遵循这一理论分析，将$\gamma$设置为1000。

简化版DCD与原始DCD在效果上基本一致，原因在于当$\gamma$足够大时，非最近邻点对的指数项贡献可以忽略不计。实验表明，简化版DCD在保持与原始DCD相近的补全质量的同时，将计算时间减少了约一个数量级，显存占用降低了约两个数量级，使得在大批量训练中使用DCD成为可能。

除了DCD之外，细粒度损失函数还包含标准CD损失和排斥损失。标准CD损失作为辅助约束，确保补全点云与真实点云之间的整体距离最小化。尽管CD存在对局部密度不敏感的缺陷，但它在衡量整体几何误差方面仍然具有不可替代的作用。CD损失的计算简单高效，能够为模型提供稳定的梯度信号，与DCD和排斥损失形成互补。

排斥损失则用于促进点云分布的均匀性，其数学定义为：
\begin{equation}
\mathcal{L}_{rep} = \frac{1}{N_s}\sum_{i=1}^{N_s}\sum_{j \in \mathcal{N}_k(\mathbf{x}_i), j \neq i} \max(0, r - \|\mathbf{x}_i - \mathbf{x}_j\|) \cdot \exp\left(-\frac{\|\mathbf{x}_i - \mathbf{x}_j\|^2}{h^2}\right)
\end{equation}

其中$N_s$为采样锚点的数量，$\mathcal{N}_k(\mathbf{x}_i)$为$\mathbf{x}_i$的$K$近邻集合，$r$为排斥半径阈值，$h$为高斯核的带宽参数。

排斥损失的设计灵感来源于PU-Net\cite{yu2018punet}中的均匀性约束。当两个预测点之间的距离小于排斥半径$r$时，排斥损失会对它们施加惩罚，惩罚的强度由高斯核权重决定。高斯核权重随着距离的增加而快速衰减，使得排斥力主要作用于非常接近的点对，而不会对距离较远的点对产生不必要的影响。这种局部排斥机制类似于物理系统中的分子间作用力，当分子距离过近时产生排斥力，防止分子过度聚集。

排斥半径$r$和带宽参数$h$的取值需要根据数据集的特性进行合理设置。对于ShapeNet归一化数据集，8192个点的平均最近邻距离约为0.02到0.03之间。排斥半径$r$应设置为平均最近邻距离的2到3倍，以确保排斥力能够有效约束局部堆积，同时不会破坏全局结构。本算法将$r$设置为0.07，$h$设置为0.03，这一参数组合在实验中表现出良好的收敛性和补全质量。

为了降低排斥损失的计算复杂度，本算法采用KNN搜索替代全量距离矩阵计算。具体而言，首先从预测点云中随机采样$N_s=2048$个锚点，然后通过KNN搜索获取每个锚点的$K=5$个最近邻点。这种采样策略将计算复杂度从$O(N^2)$降低到$O(N_s \times K)$，显著减少了计算开销和显存占用。KNN搜索通过CUDA加速的底层算子实现，进一步提升了计算效率。

细粒度损失函数由DCD损失、CD损失和排斥损失加权组合而成：
\begin{equation}
\mathcal{L}_{fine} = w_{dcd} \cdot \mathcal{L}_{dcd} + w_{cd} \cdot \text{CD}(\mathbf{P}_{fine}, \mathbf{P}_{gt}) + w_{rep} \cdot \mathcal{L}_{rep}
\end{equation}

其中$w_{dcd}$、$w_{cd}$和$w_{rep}$分别为各项损失的权重系数。

三项损失在细粒度优化中各司其职。DCD损失负责防止局部点云堆积，确保预测点与真实点之间的密度匹配；CD损失负责最小化整体几何误差，保证补全点云与真实点云之间的全局对齐；排斥损失负责促进点云分布的均匀性，避免预测点之间的过度聚集。三者的协同作用使得细粒度点云在全局结构正确的基础上，具备了更加均匀和自然的点分布特性。

(3) 损失函数协同优化

混合损失函数的有效性不仅取决于各项损失的合理设计，还依赖于多损失之间的协同优化策略。不同损失函数在训练过程中可能产生相互冲突的梯度信号，如果处理不当，会导致模型收敛困难或陷入次优解。

分阶段权重调度是本算法采用的核心协同优化策略。在训练的前50个epoch，粗粒度损失被赋予更高的权重，使得模型能够快速建立正确的全局拓扑结构。这一阶段的优化重点在于让编码器学习到有效的特征表示，并让粗粒度解码器生成合理的骨架点云。随着训练的推进，粗粒度损失的权重逐渐降低，细粒度损失的权重逐渐增加，模型的优化重心从全局结构转向局部细节。

分阶段权重调度的数学表达为：
\begin{equation}
w_{coarse}(t) = \begin{cases}
w_{coarse}^{init} \cdot \left(1 - \frac{t}{T_{warm}}\right) + w_{coarse}^{final} \cdot \frac{t}{T_{warm}}, & t < T_{warm} \\
w_{coarse}^{final}, & t \geq T_{warm}
\end{cases}
\end{equation}

其中$t$为当前epoch数，$T_{warm}$为预热阶段的epoch数（设置为50），$w_{coarse}^{init}$和$w_{coarse}^{final}$分别为粗粒度损失的初始权重和最终权重。

分阶段权重调度的设计基于对模型训练过程的深入理解。在训练初期，模型的参数随机初始化，生成的点云与真实点云之间存在较大的差距。此时如果过早地关注细粒度细节，模型可能会在错误的全局结构上浪费优化资源。通过先优化粗粒度损失，模型能够快速收敛到一个合理的全局结构，为后续的细粒度优化奠定良好的基础。这种由粗到细的优化策略与Coarse-to-Fine的生成策略形成了良好的呼应，使得模型在生成和优化两个层面都遵循了从宏观到微观的认知规律。

除了分阶段权重调度之外，本算法还采用了自动混合精度（Automatic Mixed Precision, AMP）训练策略。AMP通过在计算过程中动态选择FP16和FP32精度，在保持数值稳定性的同时显著降低了显存占用和计算时间。对于EMD和DCD等对数值精度敏感的操作，AMP会自动切换到FP32精度进行计算，确保梯度回传的准确性。AMP的引入使得本算法能够在RTX 5090 32GB显卡上以更大的批量大小进行训练，进一步提升了训练效率和模型性能。

在梯度回传方面，DCD的密度权重计算被置于\texttt{torch.no\_grad()}上下文中，避免了权重计算过程中的梯度累积，从而降低了显存占用。同时，所有距离计算和指数运算都保持在PyTorch的计算图内，确保了梯度的完整回传。这种设计在保证计算效率的同时，避免了梯度断裂导致的训练失效问题。

损失函数的协同优化还体现在各项损失之间的数值平衡上。不同损失函数的量级可能存在显著差异，如果直接相加而不进行归一化处理，量级较大的损失会主导优化过程，量级较小的损失则可能被忽略。本算法通过实验确定了各项损失的合理权重系数，使得各项损失在训练过程中能够发挥均衡的作用。权重系数的确定基于对各项损失在训练初期的数值范围的统计分析，确保各项损失对总损失的贡献大致相当。

\section{三维表面重建算法的原理与实现}

点云补全模型输出的数据本质上仍是离散的点集，难以直接应用于三维建模、渲染或3D打印等下游任务。为了构建从残缺点云到完整三维模型的系统，需要在补全后的点云基础上进行表面重建，将无序的点集转换为具有拓扑关系的网格模型。本章将详细阐述基于泊松表面重建算法的实现过程，并对重建结果进行分析与讨论。

\subsection{泊松表面重建算法原理}

(1) 泊松重建的数学基础

泊松表面重建（Poisson Surface Reconstruction）由Kazhdan等人于2006年提出，是一种基于隐式函数的表面重建方法\cite{kazhdan2006poisson}。该方法的核心思想是将表面重建问题转化为求解泊松方程的边值问题，通过点云及其法线信息构造一个指示函数，进而提取等值面作为重建的三维表面。

给定一个带有法线信息的点云集合，算法假设存在一个指示函数$\chi(\mathbf{x})$，其在模型内部取值为1，在模型外部取值为0。该函数的梯度在表面处与点云的法线方向一致，在远离表面的区域趋近于零。因此，表面重建问题可以表述为：寻找一个标量场$\chi$，使其梯度$\nabla\chi$尽可能逼近由点云法线定义的向量场$\vec{V}$。

为了实现这一目标，算法通过最小化如下能量函数来求解：
\begin{equation}
\min_{\chi} \left\| \nabla\chi - \vec{V} \right\|^2
\end{equation}

对该能量函数求导并应用散度定理，可将其转化为泊松方程：
\begin{equation}
\Delta\chi = \nabla \cdot \vec{V}
\end{equation}

其中$\Delta$为拉普拉斯算子，$\nabla \cdot \vec{V}$为向量场的散度。求解该方程后，通过提取指示函数的等值面（通常取$\chi = 0.5$）即可获得重建的三维网格表面。

与Delaunay三角化、径向基函数等传统方法相比，泊松重建在处理大规模点云时能够保持较高的重建质量和计算效率，且对点云噪声具有一定的鲁棒性\cite{kazhdan2006poisson}。

(2) 法线估计与定向

泊松重建算法的输入不仅包含点的空间坐标，还要求每个点配备正确的法线方向。法线信息的质量直接决定了重建表面的准确性，因为泊松方程的求解依赖于法线定义的向量场。

对于点云中的每个点$\mathbf{p}$，其法线估计通常基于局部邻域的几何特征。通过构建该点的$k$个最近邻点集合，计算这些点的协方差矩阵：
\begin{equation}
\mathbf{C} = \frac{1}{k}\sum_{i=1}^{k}(\mathbf{p}_i - \bar{\mathbf{p}})(\mathbf{p}_i - \bar{\mathbf{p}})^T
\end{equation}

其中$\bar{\mathbf{p}}$为邻域点的质心。对协方差矩阵进行特征分解，最小特征值对应的特征向量即为该点的法线方向。这一方法的几何直觉在于：点云的局部表面近似于一个平面，而法线方向正是该平面变化最小的方向。

然而，仅通过协方差分析得到的法线方向存在二义性——法线可能指向模型内部或外部。为了保证所有法线方向的一致性，需要对其进行定向处理。常用的方法是基于最小生成树传播法线方向，从种子点出发，使相邻点的法线夹角最小化。Open3D库提供了基于切平面一致性的法线定向算法，通过迭代优化使局部法线场保持连续且指向模型外部\cite{zhou2018open3d}。

(3) 密度阈值后处理

泊松重建过程中，算法会在整个空间域内求解泊松方程，这意味着在非表面区域也可能生成三角面片。这些冗余面片通常具有较低的密度值，因为泊松重建会为每个顶点计算一个密度属性，反映该顶点周围点云的采样密度。

密度过滤的核心思想是：表面区域的顶点应当有足够的点云支撑，而远离真实表面的顶点密度值较低。通过设定适当的密度阈值，可以有效移除这些伪影面片。具体而言，将所有顶点的密度值按升序排列，取某一低分位数（如10\%）作为阈值，将密度低于该阈值的顶点及其关联的三角面片删除。

密度阈值的选择需要在表面完整性和噪声抑制之间取得平衡。阈值过高会导致真实表面的边缘区域被误删，产生空洞；阈值过低则无法有效去除冗余面片，影响网格质量。在实际应用中，通常根据点云的采样密度和重建效果进行经验性调整。

\subsection{算法实现细节}

本章的表面重建流程基于Open3D库实现，该库提供了一套完整的点云处理与三维重建工具链\cite{zhou2018open3d}。整体实现流程涵盖点云加载、法线估计、泊松重建、密度过滤与网格平滑等关键环节。

点云数据以NumPy数组格式存储，加载后转换为Open3D的点云数据结构。本实验中使用的点云由前文所述的补全模型生成，包含8192个三维点。由于补全模型输出的点云已具备较好的几何结构，法线估计的搜索半径设置为0.1，近邻点数设置为50，以兼顾局部细节捕捉与计算效率。法线定向采用切平面一致性方法，确保所有法线指向模型外部。

泊松重建的核心参数为八叉树深度（depth），该参数决定了重建网格的分辨率。八叉树深度越大，生成的网格越精细，但计算开销也随之增加。经过实验对比，本文将depth设置为9，该参数配置下重建的网格包含超过20000条边和40000个三角面片，在细节保留与计算效率之间取得了较好的平衡。重建尺度参数（scale）设置为1.1，略微放大的重建范围可以防止模型边界处的几何特征被裁剪。

重建完成后，密度过滤采用10\%分位数作为阈值，移除低密度区域的冗余面片。随后执行退化三角形去除、重复顶点合并等拓扑清理操作。为进一步提升表面质量，本文采用Taubin平滑算法对网格进行处理。Taubin平滑通过交替使用正负拉普拉斯流，在平滑表面的同时避免网格收缩\cite{taubin1995signal}。平滑迭代次数设置为5次，该参数能够在消除高频噪声的同时保留模型的几何细节。最终重建的网格被保存为PLY格式，并统一着色为浅灰色以便于可视化。


\section{本章小结}

本章系统阐述了基于Transformer的点云补全算法与基于泊松方程的三维表面重建算法的整体架构设计与实现细节。在点云补全方面，算法采用Transformer经典的编码器-解码器的基本框架，编码器通过点云序列化和几何感知自注意力机制提取输入点云的全局特征，解码器通过Coarse-to-Fine渐进式生成策略逐步生成完整点云。针对补全过程中出现的点云密度不均匀问题，本文设计了包含EMD损失、密度感知倒角距离和排斥损失的混合损失函数，通过分阶段权重调度策略实现多损失协同优化，有效提升了补全点云的全局结构一致性和局部分布均匀性。在三维表面重建方面，本文构建了从补全点云到三维网格的完整重建流程，包括法线估计、泊松重建、密度过滤和Taubin平滑等关键步骤。该重建流程能够基于补全后的高质量点云生成表面连续、细节丰富的三角网格模型，为点云补全算法的实际应用提供了有效的下游任务验证。
